\section{Introduction}

\subsection{Background and Motivation}
Modern smart city initiatives rely heavily on distributed Internet of Things (IoT) sensor grids to maintain situational awareness across vital public utilities. Continuous telemetry streams generated by municipal sensors monitor pipeline flow and hydrostatic pressure in potable water distribution networks, load and power factor across electrical distribution grids, and particulate matter ($\text{PM}_{2.5}$, $\text{PM}_{10}$) and hazardous gases in municipal air quality networks. These telemetry feeds inform automated utility dispatch, preventive maintenance algorithms, and public health advisories.

However, deploying thousands of low-cost virtual and physical sensing nodes in public and semi-public environments introduces severe cybersecurity and operational challenges. Edge sensors typically operate over unencrypted or minimally protected wireless networks, making them susceptible to spoofing, message alteration, and hostile denial-of-service (DoS) bursts. Furthermore, sensor transmissions frequently embed precise spatial coordinates (such as GPS latitude and longitude) and hardware identifiers (such as network interface MAC addresses) to aid in sensor discovery and provisioning. When stored unredacted in public or municipal data lakes, these persistent identifiers create severe privacy risks for citizens whose residences or workplaces correlate with sensor installation sites.

\subsection{Problem Statement}
Conventional IoT data pipelines generally route incoming sensor data directly from an ingestion transport (such as an MQTT broker or HTTP endpoint) into cloud time-series databases, deferring validation and data cleaning to downstream batch analytics jobs. This conventional design exhibits critical vulnerabilities:
\begin{enumerate}[label=(\alph*)]
    \item \textbf{Ingestion Poisoning and Cascading Failures:} Malicious or broken sensors can transmit rapid bursts of invalid or physically impossible readings (e.g., negative water consumption or out-of-scale air quality indices). Storing these corrupted readings taints municipal monitoring algorithms and consumes database I/O capacity.
    \item \textbf{Privacy Infiltration:} Once telemetry records containing raw geographic coordinates and device MAC addresses are written to relational or time-series storage, retroactive sanitization is computationally expensive and legally hazardous under data privacy regulations such as the General Data Protection Regulation (GDPR) and municipal data privacy mandates.
    \item \textbf{Cascade Resource Exhaustion:} If the underlying database suffers latency degradation or transient unavailability, unbuffered synchronous ingestion endpoints quickly exhaust thread pools and socket descriptors, causing the entire gateway to crash.
\end{enumerate}

\subsection{Project Objectives}
The primary objective of the \textbf{Secure IoT Gateway Simulation for Resource Sensing (SIGS-RS)} project is to engineer an enterprise-grade, asynchronous edge gateway engine that guarantees that unverified or unsanitized telemetry \emph{never reaches persistent database storage}. Specifically, the project addresses the following engineering objectives:
\begin{itemize}
    \item \textbf{Multi-Transport Secure Ingestion:} Support dual-transport ingestion via HTTP REST and MQTT over TLS 1.2+ (MQTTS on port 8883) with hard size limits ($\le 512$~KB per packet) to protect edge gateway memory.
    \item \textbf{Sequential Two-Stage Middleware Pipeline:} Enforce a strict sequential invariant where every packet is inspected for rate anomalies and physical impossibilities (Stage 1) before proceeding to privacy anonymization (Stage 2). Packets rejected in Stage 1 are dropped immediately and logged to an immutable audit trail without entering Stage 2.
    \item \textbf{Fail-Secure Privacy Anonymization:} Automatically strip all spatial coordinates (\texttt{latitude}, \texttt{longitude}, \texttt{geo}) and device identifiers (\texttt{mac\_address}, \texttt{device\_id}) prior to persistence.
    \item \textbf{Resilience and Fail-Closed Operation:} Integrate an adaptive circuit breaker to protect against database write latency degradation ($> 2,000$~ms) and enforce automated simulation lockout upon compliance guideline revisions (SRS BR-3).
    \item \textbf{Empirical Verification and Traceability:} Validate all software requirements through automated testing (580 unique test cases) and reproducible microbenchmarking suites.
\end{itemize}

\subsection{Scope and Delimitations}
The scope of this project encompasses the backend simulation engine, asynchronous ingestion layer, two-stage security middleware, database repository, circuit breaker, management API, virtual sensor load simulator, and containerization manifests. 

To preserve technical honesty, a clear delimitation is established between the \emph{locally verified single-process development environment} and \emph{production cluster prerequisites}. High-scale throughput targets (50,000 packets/second across 10,000 concurrent virtual sensors per REQ-PF-1) and hardware-enforced Write-Once-Read-Many (WORM) storage represent architectural cluster goals that require multi-node staging infrastructure and are explicitly designated as unverified on the local host.

\subsection{Report Organization}
The remainder of this report is structured as follows: Section~2 reviews related literature and consulted technical standards. Section~3 examines the limitations of existing IoT systems. Section~4 presents the proposed SIGS-RS system. Section~5 provides a comprehensive synthesis of the Software Requirements Specification (SRS) and traceability matrix. Section~6 details the system design and architectural models. Section~7 reviews the confirmed technology stack. Sections~8 and~9 present working modules and detailed component descriptions. Section~10 discusses the test plan and empirical test outcomes. Section~11 evaluates performance benchmarking results. Section~12 presents screenshots and execution evidence. Section~13 concludes the report, and Section~14 discusses future work. Implementation artifacts are included in Appendix~A, followed by the verbatim SRS in Appendix~B.
